\chapter{اللقاء الثالث عشر: قصة أصحاب السبت ومسخهم قردة}

\section{تأويل (ولقد علمتم الذين اعتدوا منكم في السبت)}
السلام عليكم ورحمة الله وبركاته. الحمد لله رب العالمين، والصلاة والسلام على نبينا محمد وعلى آله وصحبه أجمعين.
وصلنا بحمد الله تبارك وتعالى إلى المجلد الثاني من تفسير الطبري في قوله جل ثناؤه: \begin{ayah}وَلَقَدْ عَلِمْتُمُ الَّذِينَ اعْتَدَوْا مِنكُمْ فِي السَّبْتِ فَقُلْنَا لَهُمْ كُونُوا قِرَدَةً خَاسِئِينَ\end{ayah}.

\isnad{قال الطبري رحمه الله:} «يعني بقوله \textit{ولقد علمتم}: ولقد عرفتم، كقولك: قد علمت أخاك ولم أكن أعلمه، يعني عرفته ولم أكن أعرفه، كما قال جل ثناؤه: \begin{ayah}وَآخَرِينَ مِن دُونِهِمْ لَا تَعْلَمُونَهُمُ اللَّهُ يَعْلَمُهُمْ\end{ayah} يعني لا تعرفونهم الله يعرفهم».

\begin{fawaid}[الفرق بين العلم والمعرفة في الاستعمال القرآني]
المعرفة أقل من العلم؛ فالمعرفة تأتي بمعنى الإدراك المجرد الذي قد لا يتبعه عمل، كما قال تعالى: \textit{يعرفونه كما يعرفون أبناءهم}، \textit{يعرفون نعمة الله ثم ينكرونها}. أما العلم في القرآن والسنة فكثيراً ما يتضمن معنى الإدراك المتبوع بالخشية والعمل والاتباع، كقوله ﷺ: \textit{أنا أعلمكم بالله وأشدكم له خشية}، وقوله تعالى: \textit{إنما يخشى الله من عباده العلماء}. 
وتفسير الطبري للعلم هنا بالمعرفة له وجه قوي في هذا السياق؛ لأن المراد: أدركتم وبلغكم خبر هؤلاء.
\end{fawaid}

\isnad{قال الطبري:} «وقوله \textit{الذين اعتدوا منكم في السبت} أي الذين تجاوزوا حدي وركبوا ما نهيتهم عنه في يوم السبت وعصوا أمري... وحذر المخاطبين بها أن يحل بهم بإصرارهم على كفرهم ومقامهم على جحود نبوة محمد ﷺ مثل الذي حل بأوائلهم من المسخ والرجف».

\separator

\section{العبرة من القصة والابتلاء بقرب المعصية}
ساق \rawi{الطبري} الآثار والإسرائيليات في تفصيل قصة أصحاب السبت، وكيف أن الحيتان كانت تشرع لهم يوم السبت (تظهر على الماء)، وتختفي في غيره، فاحتالوا باصطيادها وحبسها يوم السبت وأخذها يوم الأحد.

\begin{fawaid}[سنة الابتلاء بقرب المعصية وسهولتها]
من سنن الله في خلقه أن يبتلي العبد بتيسير أسباب المعصية وقربها من متناول يده، كما ابتلى أصحاب السبت بظهور الحيتان، وكما ابتلى الصحابة في الإحرام بالصيد، قال تعالى: \textit{ليبلونكم الله بشيء من الصيد تناله أيديكم ورماحكم ليعلم الله من يخافه بالغيب}. وهذا من أعظم الامتحانات في زماننا مع تيسر الوصول للمعاصي عبر الأجهزة الذكية وغيرها.
\end{fawaid}

\begin{fawaid}[أصناف الناس عند ظهور المنكر]
انقسم الناس في هذه القصة إلى ثلاثة أصناف:
\begin{itemize}
    \item \textbf{صنف اتقى ونهى:} وهم الذين استمسكوا بالحق وتواصوا به ونهوا عن المنكر، وهؤلاء هم الربانيون، وهم أعلى الناس درجة.
    \item \textbf{صنف أمسك ولم ينهَ:} وهم الذين لم يعصوا، ولكنهم ضعفوا عن إنكار المنكر وتثبيط العصاة.
    \item \textbf{صنف اعتدى وعصى:} وهم الذين انتهكوا حرمة الله، فنزل بهم العذاب والمسخ.
\end{itemize}
\end{fawaid}

\separator

\section{تأويل (فقلنا لهم كونوا قردة خاسئين) والرد على مجاهد}
\isnad{قال أبو جعفر:} «يعني بقوله \textit{فقلنا لهم} أي فقلنا للذين اعتدوا في السبت... \textit{كونوا قردة خاسئين} أي صيروا كذلك، والخاسئ: المبعد المطرود، كما يُخسأ الكلب... أي مبعدين من الخير، أذلاء صغراء».

\subsection{رد الطبري على من أوّل المسخ}
أورد \rawi{الطبري} أثراً عن \rawi{مجاهد} في قوله \textit{كونوا قردة خاسئين} أنه قال: «مسخت قلوبهم ولم يمسخوا قردة، وإنما هو مثل ضربه الله لهم كمثل الحمار يحمل أسفاراً».

فرد عليه \rawi{الطبري} بقوة قائلاً: «وهذا القول الذي قاله \rawi{مجاهد} قولٌ لظاهر ما دل عليه كتاب الله مخالف... فسواء قال قائل: هم لم يمسخهم قردة، وقد أخبر جل ذكره أنه جعل منهم قردة وخنازير، وآخر قال: لم يكن شيء مما أخبر الله عن بني إسرائيل أنه كان منهم... ومن أنكر شيئاً من ذلك وأقر بآخر منه، سُئل البرهان على قوله... ثم يُسأل الفرق من خبر مستفيض أو أثر صحيح، هذا مع خلاف قول \rawi{مجاهد} قول جميع الحجة التي لا يجوز عليها الخطأ والكذب فيما نقلته مجمعة عليه، وكفى دليلاً على فساد قولٍ إجماعُها على تخطئته».

\begin{fawaid}[قاعدة الطبري الصارمة في رد الأقوال الشاذة]
أصّل \rawi{الطبري} هنا قاعدة متينة في رد الأقوال الشاذة وإن صدرت من إمام كبير كـ \rawi{مجاهد}؛ وهي: لا يجوز ترك المفهوم من ظاهر التنزيل إلى معنى باطن لا دلالة عليه من ظاهر القرآن، ولا من خبر صحيح عن الرسول ﷺ، ولا أثر مستفيض، مع مخالفته لإجماع المفسرين (الحجة). فإجماعهم على خلافه يكفي دليلاً على فساد القول.
\end{fawaid}

\separator

\section{تأويل (فجعلناها نكالاً لما بين يديها وما خلفها)}
اختلف المفسرون في الضمير (الهاء والألف) في قوله \textit{فجعلناها} على ماذا يعود؟
ذكر \rawi{الطبري} الأقوال:
\begin{itemize}
    \item روي عن \rawi{ابن عباس} أنها تعود على العقوبة (المَسْخة).
    \item وروي عنه أيضاً أنها تعود على الحيتان.
    \item وقال آخرون تعود على القرية التي اعتدى أهلها.
    \item وقال آخرون تعود على القردة.
    \item وقال آخرون تعود على الأمة التي اعتدت.
\end{itemize}

\begin{fawaid}[منهج الطبري في تأخير الترجيح واستيعاب الاحتمالات]
من منهج \rawi{الطبري} أنه يورد الأقوال والاحتمالات في عودة الضمير، ثم يشرح معنى الآية بناءً على كل احتمال، ثم يؤخر ترجيحه للقول الصحيح لصفحات عديدة. وهنا رجح بعد أربع صفحات أن الضمير يعود على "العقوبة والمَسْخة"، وضعّف عودته على "الحيتان" لأنه أبعد في السياق ولم يجرِ لها ذكر صريح.
\end{fawaid}

\isnad{وأما تأويل النكال:} فقال \rawi{الطبري}: «والنكال مصدر من قول القائل: نكل فلان بفلان تنكيلاً ونكالاً، وأصل النكال العقوبة».

\separator

\section{تأويل (وموعظة للمتقين) والانتفاع بالقرآن}
\isnad{قال أبو جعفر:} «والموعظة مصدر من قول القائل: وعظت الرجل أعظه وعظاً وموعظة... ليتعظوا بها ويعتبروا ويتذكروا بها... فجعل تعالى ذكره ما أحل بالذين اعتدوا في السبت من عقوبته موعظة للمتقين خاصة، وعبرة للمؤمنين دون الكافرين به إلى يوم القيامة».

\begin{fawaid}[قصص القرآن للعمل لا للتسلية]
قصص القرآن لم تُسق لمجرد التوثيق التاريخي أو التسلية، بل للموعظة والاعتبار. ولكن هذه العبرة لا ينتفع بها إلا "المتقون" الذين صلحت قلوبهم وطابت، كالأرض الطيبة التي تقبل الماء فتنبت الكلأ. فإذا تُليت عليك أخبار القوم ولم يتحرك قلبك للعظة والعمل، فاعلم أن في القلب نقصاً ومرضاً يحتاج إلى معالجة.
\end{fawaid}

\separator

\section{تأويل (وما الله بغافل عما تعملون)}
\isnad{قال أبو جعفر:} «وما الله بغافل يا معشر المكذبين بآياته... عما تعملون من أعمالكم الخبيثة... وأصل الغفلة عن الشيء: تركه على وجه السهو عنه والنسيان له. فأخبرهم تعالى ذكره أنه غير غافل عن أفعالهم الخبيثة ولا ساهٍ عنها، بل هو لها محصٍ ولها حافظ».

\begin{fawaid}[إثبات الكمال لله بنفي النقص]
الصفات والأفعال المنفية عن الله تبارك وتعالى (كالغفلة، والظلم، واللغوب، والسنة، والنوم) لا تُنفى لمجرد النفي، بل لأن نفي النقص يدل على إثبات كمال الضد. فنفي الظلم يثبت كمال العدل، ونفي الغفلة يثبت كمال العلم والمراقبة والإحاطة، ونفي اللغوب (التعب) يثبت كمال القدرة.
\end{fawaid}

\separator

\section{نصيحة عامة في تزكية النفس وإصلاح القلب}
العلم والمعرفة إذا وردا على قلب سليم انتفع بهما العبد وزكّيا روحه، وإذا وردا على قلب قاسٍ ممتلئ بالفتن لم يزداه إلا كبراً وطغياناً، كما قال تعالى في المنافقين: \textit{فزادتهم رجساً إلى رجسهم}. 
لذلك يجب على طالب العلم أن يجعل غايته من دراسة التفسير والسنن هو أن يزكي نفسه أولاً، وأن يُقوّم أخلاقه، وأن يجعل القرآن دليلاً له يتحرك به في الناس، لا أن يكون مجرد جامع للمعلومات ليجادل بها في المجالس ووسائل التواصل. 

قال الإمام \rawi{ابن تيمية} رحمه الله في الوصية الصغرى ما معناه: (فمن نوّر الله قلبه هداه بما بلغه من العلم، ومن لم يكن كذلك فلن تزيده كثرة الكتب والعلوم إلا حيرة وضلالاً).

بهذا نكون قد مررنا على أهم فوائد هذا الدرس، ونكمل في اللقاء القادم بإذن الله. والسلام عليكم ورحمة الله وبركاته.